\section{Purpose, notation, and the assertion still to be proved}
\label{sec:purpose}

This manuscript begins a conventional reconstruction of the announced proof
of the Riemann hypothesis and its extension to Dirichlet $L$-functions
\cite{alpo-announcement}.
It is a mathematical working paper, not a completed proof of those
hypotheses. The results below have proofs in ordinary mathematical
language. Where a result from the supplied papers is used as an input,
that input is stated explicitly. In particular, the arithmetic estimate
\eqref{eq:remaining-estimate} is not proved in this version. Calling it
an auxiliary lemma would not make the reconstruction complete.

The aim is to separate the analytic conclusion, which can be proved quite
concisely, from the arithmetic construction that must produce the required
estimate. This also identifies precisely what can be gained from the
earlier work on zero-free regions and from the quadratic determinant
arguments. The latter arguments are reconstructed below because their
divisibility and nonvanishing mechanisms are useful in their own right;
their conclusions do not supply \eqref{eq:remaining-estimate}.

We write $s=\sigma+it$, use natural logarithms, and put
$y_+=\max(y,0)$. The notation $A\ll B$ means $|A|\le CB$ for a
positive constant; subscripts specify permitted dependencies.
A Dirichlet character modulo $q$ is a homomorphism
$(\Z/q\Z)^\times\to\C^\times$, extended periodically to $\Z$ by
zero on integers not coprime to $q$. It is primitive if it is not induced
from a proper divisor of $q$. Its conductor is the least modulus from
which it is induced. The character of conductor one is denoted by $1$.
For $\sigma>1$,
\[
 L(s,\chi)=\sum_{n\ge1}\frac{\chi(n)}{n^s}
 =\prod_p(1-\chi(p)p^{-s})^{-1},\qquad L(s,1)=\zeta(s).
\]
The von Mangoldt function is $\Lambda(p^j)=\log p$ for primes $p$
and integers $j\ge1$, and $\Lambda(n)=0$ otherwise.

We use the classical analytic continuation and functional equations of
these functions, and their nonvanishing on $\Re s=1$ away from the
pole of $\zeta$. In particular $L(1,\chi)\ne0$ for nonprincipal
$\chi$. For primitive $\chi$ of conductor $q>1$, with
$\chi(-1)=(-1)^a$, the completed function
\[
 \Lambda_\chi(s)=(q/\pi)^{(s+a)/2}
       \Gamma((s+a)/2)L(s,\chi)
\]
is entire and satisfies
$\Lambda_\chi(s)=w_\chi\Lambda_{\bar\chi}(1-s)$ with
$|w_\chi|=1$. Also
$\zeta(s)=(s-1)^{-1}+\gamma+O(s-1)$ at $s=1$.
These classical facts are background inputs; the arguments below do not
claim to reprove the analytic theory of $L$-functions from its
foundations. References and normalizations are given in
\cite{dlmf-dirichlet,dlmf-zeta,dlmf-zeros}.

For $x\ge1$ define the smoothed sums
\begin{align}
 S_\chi(x)&=\sum_{n\le x}\Lambda(n)\chi(n)(x/n-1),\label{eq:Schi}\\
 G_\chi(x)&=S_\chi(x)+x\frac{L'(1,\chi)}{L(1,\chi)}
                  &&(\chi\ne1),\label{eq:Gchi}\\
 G_1(x)&=\sum_{n\le x}\Lambda(n)(x/n-1)
          -x\log x+(1+\gamma)x.\label{eq:Gone}
\end{align}
Only primitive characters will be used in these definitions. In
particular, \eqref{eq:Gone} refers to conductor one, not to the
principal character at an arbitrary larger modulus.

Fix a primitive character $\chi$ modulo $q$, let $o$ be its order, and
let $\mathcal C_\chi$ be the set of distinct primitive characters
inducing the powers $\chi^j$, $0\le j<o$. Put
\begin{equation}\label{eq:Hchi}
 H_\chi(x)=G_1(x)+\frac1{3o+1}
       \sum_{\substack{\psi\in\mathcal C_\chi\\\psi\ne1}}
               \Re G_\psi(x).
\end{equation}
The set and its weights are fixed once $\chi$ is fixed; they do not
change with $x$ or with an error exponent.

Here is a precise sufficient arithmetic assertion. For each primitive
$\chi$, there is a fixed $b>1$ such that, for every $\varepsilon>0$,
there are arbitrarily large $Y$ for which
\begin{equation}\label{eq:remaining-estimate}
             H_\chi(x)\le x^{1/2+\varepsilon}
                    \qquad(Y\le x\le Y^b).
\end{equation}
The proof below shows that this assertion implies Dirichlet GRH.
Its proof from the proposed automorphic and arithmetic constructions
remains the principal unfinished part of the present manuscript.
The one-sided inequality and the quantifier ``arbitrarily large'' are
intentional: an absolute-value bound at every sufficiently large $x$
is unnecessary.

There is a weaker sufficient formulation.  For the same fixed function
$H_\chi$, it suffices to find $\delta_j\to0$, $M_j\geq0$ and $b_j>1$
such that
\[
 H_\chi(x)\leq M_jx^{1/2+\delta_j}
       \qquad(Y\leq x\leq Y^{b_j})
\]
on arbitrarily late intervals for each fixed $j$.
Theorem~\ref{thm:vanishing-loss-windows} proves this assertion.
The constants and interval shapes may depend on $j$; the target
function does not.  Thus a construction need only make its total
exponent loss tend to zero, while controlling every $x$ in the
intervals belonging to each fixed accuracy. More generally,
Corollary~\ref{cor:varying-positive-families} permits a new finite positive
combination $H_j$ at each accuracy, provided $\chi$ or $\bar\chi$ occurs
with a positive coefficient. Each $H_j$ remains fixed during its own
sequence of intervals. A common auxiliary family across accuracies is
therefore unnecessary.

The arithmetic programme has several distinct requirements.
Theorem~\ref{thm-classical-galois-zero-monodromy} supplies the classical
Galois representations and their required local and infinity data.
Integral theta maps must exist and be nonzero, with their
coefficient lattices, periods and divisibility fixed precisely.
Their global integral identities and analytic norms must retain the
same measures and normalizations. The geometric argument also needs
a bound for normalized logarithmic covolumes of lattices associated
with Hodge bundles, the bundles of invariant differentials in the
relevant families of abelian varieties. Finally, the cohomological,
spectral and geometric estimates must apply to these same objects.
Neither local residue positivity nor an abstract lattice comparison
alone supplies these global constructions.

One may aim to construct the required compatible data separately
at each scale. A larger family construction is a sufficient method,
but its auxiliary organization is not part of the analytic conclusion.
Replacing it requires an actual construction at every scale in the
chosen intervals, with all of its estimates and compatibility conditions.
Proposition~\ref{cs-finite-average} gives a more flexible arithmetic
target: an averaged discrepancy and explicit finite errors.
Corollary~\ref{cs-exceptional-budget} likewise retains the actual
discarded-degree and index loss in the covolume comparison.

